\chapter{Complete Technology Stack}

\section{Ecosystem Overview}
The system architecture combines a Python backend runtime with a modern JavaScript/React web client. Every dependency included in the repository has been audited and cataloged with its version, architectural role, trade-offs, and typical interview questions.

\section{Backend and Machine Learning Technologies}

\subsection{Python 3.8+ (Core Runtime)}
\begin{itemize}[leftmargin=*]
    \item \textbf{Role}: Underlying runtime for all mathematical calculations, computer-vision pipelines, and Flask REST web servers.
    \item \textbf{Why Used}: Unrivaled ecosystem for scientific computing, deep learning integration (PyTorch, Ultralytics), and high developer velocity.
    \item \textbf{Interview Focus}: \textit{How does Python's Global Interpreter Lock (GIL) affect multi-threaded computer-vision inference?}
    \item \textbf{Spoken Explanation}: ``In \texttt{backend/app.py}, the CV inference loop runs in a background daemon thread while Flask serves requests. Because OpenCV and PyTorch release the GIL during native C++ tensor and matrix computations, the Flask HTTP worker can handle incoming requests without being blocked by the heavy visual pipeline.''
\end{itemize}

\subsection{Ultralytics YOLOv8n (Object Detection)}
\begin{itemize}[leftmargin=*]
    \item \textbf{Role}: Deep convolutional neural network for single-pass object detection across road vehicles.
    \item \textbf{Why YOLOv8n}: The nano variant (\texttt{yolov8n.pt}, $\approx 6.2\,\text{MB}$) provides the optimal trade-off between inference speed ($\approx 30\,\text{ms}$ on modern CPUs) and mean Average Precision (mAP) for edge applications without dedicated server GPUs.
    \item \textbf{Trade-offs}: Nano weights sacrifice fine-grained detection of partially occluded distant vehicles compared to larger variants (\texttt{yolov8x.pt}), but avoid CPU freezing.
\end{itemize}

\subsection{OpenCV (Computer Vision Utilities)}
\begin{itemize}[leftmargin=*]
    \item \textbf{Role}: Video stream ingestion, frame decoding, coordinate drawing, visual HUD overlay generation, and spatial transformations (\texttt{cv2.line}, \texttt{cv2.putText}).
    \item \textbf{Why Used}: High-performance C++ native routines with minimal overhead for frame-by-frame video manipulation.
\end{itemize}

\subsection{scikit-learn 1.9.0 (Random Forest Classifier)}
\begin{itemize}[leftmargin=*]
    \item \textbf{Role}: Supervised machine-learning model in \texttt{backend/congestion\_predictor.py} classifying vehicular density into congestion tiers.
    \item \textbf{Why Used}: Random Forest ensembles mitigate overfitting on non-linear decision boundaries through bagging and decision tree majority voting.
\end{itemize}

\subsection{Flask 3.1.3 and Flask-CORS 6.0.5 (REST API Layer)}
\begin{itemize}[leftmargin=*]
    \item \textbf{Role}: Exposing HTTP REST endpoints for state retrieval, signal queries, emergency simulation, and commuter telemetry.
    \item \textbf{Why Used}: Extremely lightweight micro-framework avoiding the monolithic overhead of Django, easily deployable in memory-constrained cloud containers.
\end{itemize}

\section{Frontend Technologies}

\subsection{React 18.3.1 and Vite 5.2.11}
\begin{itemize}[leftmargin=*]
    \item \textbf{Role}: Declarative component-driven user interface and modern ES-module build bundling.
    \item \textbf{Why Used}: React's virtual DOM reconciliation efficiently updates real-time animated dashboards without full-page reloads. Vite provides instantaneous Hot Module Replacement (HMR) and optimized Rollup tree-shaking production builds.
\end{itemize}

\subsection{Zustand 4.5.2 (State Management)}
\begin{itemize}[leftmargin=*]
    \item \textbf{Role}: Centralized reactive store managing traffic telemetry, signal configurations, historical buffers, and commuter advice state.
    \item \textbf{Why Zustand over Redux}: Provides a clean hooks-based API without boilerplate reducers, actions, or Context provider re-render cascades.
\end{itemize}

\subsection{Tailwind CSS 3.4.4 and Framer Motion 11.2.10}
\begin{itemize}[leftmargin=*]
    \item \textbf{Role}: Custom cyberpunk visual identity, glassmorphic card panels, ambient glow accents, and spring-based hardware-accelerated transitions.
\end{itemize}

\subsection{Recharts 2.12.7}
\begin{itemize}[leftmargin=*]
    \item \textbf{Role}: Declarative SVG rendering of time-series vehicular density streams, stacked violation bar charts, and radar charts.
\end{itemize}

\section{Infrastructure and Storage}
\begin{itemize}[leftmargin=*]
    \item \textbf{MongoDB / PyMongo 4.17.0}: Document-oriented database for persistence of infraction snapshots and hourly traffic aggregate statistics, implemented with a graceful fallback to standard output when offline.
    \item \textbf{Vercel}: Edge CDN hosting the React SPA with automated SPA routing rewrite rules (\texttt{frontend/vercel.json}).
    \item \textbf{Render}: Containerized cloud hosting platform running the Python web service (\texttt{backend/app\_demo.py}).
\end{itemize}
